\clearpage\fancyhead[L]{\small COMET.}\SourceLeaf{350}
\begin{center}\textbf{COMET.}\par\vspace{4pt}\hrule\vspace{4pt}\fontsize{10}{12}\selectfont Monatliche Correspondenz. Vol. XXV, pp. 94..97, January 1812.\par\vspace{4pt}\hrule\vspace{4pt}\end{center}

\par\indent [The contents of these letters, apart from the last paragraph, are also given in the article of 25 January 1812 in the Gött. gel. Anz. Here, however, the time of perihelion is given as .\;.\;.\;.\;. 1811 Nov. 11,392108, Göttingen meridian.]

\par\indent Mr. \textsc{Wachter} has also recently calculated the occultation of $\lambda$ Aquarii and found the following improved conjunction times.

\begin{center}
\fontsize{10}{12}\selectfont\setlength{\tabcolsep}{6pt}\renewcommand{\arraystretch}{1.15}\begin{tabular}{l|ccc}
Göttingen & $10^{\text{u}}$ & $22^{\text{m}}$ & $29^{\text{s}}69$ \\
Mannheim & 10 & 16 & 38.01 \\
Cronach & 10 & 28 & 3.50 \\
\end{tabular}
\end{center}

\noindent The first two agree well with yours, but the last differs by 10 seconds. Since he did the calculation twice, using two different methods, you may find it worth checking your last calculation once more.
